\documentclass[ecta,nameyear,draft]{econsocart}
\usepackage{amsmath,amsfonts,amssymb,amsthm,booktabs,array,longtable}
\usepackage[T1]{fontenc}
\usepackage{mathpazo,microtype,needspace,xurl}
\RequirePackage[colorlinks,citecolor=blue,linkcolor=blue,urlcolor=blue]{hyperref}
\startlocaldefs
\newtheorem{theorem}{Theorem}
\newtheorem{proposition}[theorem]{Proposition}
\newtheorem{lemma}[theorem]{Lemma}
\newtheorem{corollary}[theorem]{Corollary}
\newcommand{\E}{\mathbb E}
\newcommand{\R}{\mathbb R}
\endlocaldefs
\setlength{\emergencystretch}{3em}
\makeatletter\def\copyright@text{Prepared for referee review}\makeatother

\setcounter{secnumdepth}{2}

\begin{document}
\begin{frontmatter}
\title{Neural Bellman Operators}
\runtitle{Neural Bellman Operators}
\begin{aug}
\author[id=au1,addressref={add1}]{\fnms{Qian}~\snm{QI}\ead[label=e1]{qiqian@pku.edu.cn}}
\address[id=add1]{Peking University}
\end{aug}
\begin{abstract}
This paper develops Neural Bellman Operators as constructive policy-evaluation and feasible-improvement procedures. An own-future neural continuation retains its action witnesses and admits an explicit policy-loss account. For the original nonlinear investment economy, primitive convexity and semiconcavity support a second implementation: recursive Bellman queries replace the state lattice, while a trained selector routes a query only after its accuracy certificate remains unresolved. We prove uniform original-optimum guarantees for the complete acquired policy, finite safeguarded refinement, and an explicit state-dimension and horizon work account. A certificate-localization result identifies when an additional prediction is sufficient to terminate verification. Two source-frozen catalogues contain 256 complete-process services in two, eight and sixteen state dimensions. Gating reduces the deployment queries of unconditional prediction insertion, but learned services do not beat the conventional controls in cold-start process time on these workloads. Earlier continuous-law comparisons in which pure fitted policies are worse remain separate evidence. The economic purchase criterion distinguishes policy accuracy, expected gain, and the price of computation.
\end{abstract}
\begin{keyword}
\kwd{Dynamic programming}\kwd{Neural approximation}\kwd{Policy certification}
\kwd{Information acquisition}\kwd{Computational economics}
\end{keyword}
\end{frontmatter}
\section{Introduction}\label{sec:introduction65}
A continuation approximation is useful because it helps produce an economic decision. Small fitting error, a verified bound on the loss of the implemented policy, and inexpensive computation are different achievements. A method that improves one need not improve the others. This paper connects these objects within a constructive Neural Bellman Operator (NBO), and studies which role learning can play when the decision must satisfy a stated accuracy contract.

The first construction forms each continuation against its own fitted future, retains feasible action witnesses, and accounts for state coverage, action search, integration and numerical implementation. Its exact affine--ReLU realization does not require a successful nonconvex training run. A second, trained implementation uses a network to propose an action query. It verifies that query through the original Bellman problem rather than taking the network's output as a value certificate. The latter is an action-selector implementation, not a relabeling of an inaccurate fitted-value policy.

Our original investment economy has nonlinear transitions, a state-dependent capacity constraint, and continuous innovations. Primitive state curvature dominates the transition's convexity defect. The resulting Lipschitz and semiconcavity bounds permit conditional-mean quadrature and lower bounds on the full continuous action interval. A recursive query obtains a continuation only at the descendants it actually needs. It does not construct a conventional policy at every point of a state lattice before allowing the neural proposal to act.

The key computational distinction is between the need for another query and its proposed location. Both boundary actions are first evaluated. If their certificate already meets the requested accuracy, no prediction is evaluated. Otherwise a learned proposal can replace a required safeguarded split inside an unresolved interval. Its numerical content has no authority over the lower bound, feasibility test, or stopping tolerance. A conventional parabolic rule supplies the comparison, and an explicitly costed rational recovery handles a failed floating certificate without consulting a stored reference policy.

Theorem~\ref{thm:query63} connects the resulting query to the complete policy's own future under the original continuous law. Theorem~\ref{thm:gated64} gives its safeguarded termination and prediction-free return properties. Proposition~\ref{prop:localization65} further separates two requirements for a useful learned query: it must supply a good witness and leave a sufficiently localized certificate. These statements concern implementable policies and effective verification, not a presumption that training is successful or that a neural representation is uniquely capable of certification.

The experiments retain that distinction. The tensor-reference comparison produced twenty pure ReLU policies with strictly higher expected cost than the common policy, while augmented comparisons were unresolved and the common method had the smallest complete work in every attainable task--target cell. We preserve those results. The subsequent recursive-query catalogues instead test a different computational organization on the same scalar investment primitives. They include conventional adaptive parabolic search, bisection, ReLU and quadratic proposals, and a matched insertion-versus-gating comparison. Every label, fit, query and durable output is charged to its isolated process. The two catalogues contain 96 and 160 services, respectively; their deterministic workloads are not additional independent continuous-law cost observations.

This produces a positive state-storage result without concealing the comparative result. Callable policies attain the declared contracts without a stored state lattice, including at sixteen state coordinates. Gating removes avoidable prediction work and sometimes reduces the number of Bellman queries relative to adaptive conventional search. Nevertheless, the learned services remain more expensive at cold start on the deposited workload. The value of a fitted query, its acquisition cost, and the horizon dependence of verification must therefore be examined together. The work account makes these quantities observable rather than substituting a successful fallback or a favorable online count for a full service comparison.

\paragraph{Relation to existing methods.}
Lipschitz extension and distance transforms underlie the constructive cone and its compiler \citep{mcshane1934,felzenszwalb2012}. Approximate policy iteration and safe policy improvement have established error and performance analyses \citep{munos2003,pirotta2013,scherrer2015}. Convex Bellman approximation is studied by \citet{yang2020convex}. Here primitive curvature is established for the specified bilinear capacity-constrained economy and is carried through actual acquired actions, continuous-action search and the policy's own future. The proof does not assume that an arbitrary fitted neural continuation is convex.

Learning a useful optimization initialization is also not a new principle. \citet{sakaue2023warm} develop prediction-dependent warm starts for discrete-convex minimization, and \citet{agrawal2020control} learn parameters of optimization-based control policies. Our setting instead requires outward original-law Bellman intervals at recursively queried states. The additional certificate-demand rule delays a prediction until the current information is insufficient; it leaves the correctness argument independent of prediction quality. The comparisons identify the work of that rule separately from the common analytic terminal kernel and from the cost of fitting.

\paragraph{Organization and scope.}
Sections~\ref{sec:core49} and \ref{sec:construction49} retain the original operator, own-future construction and witness-preserving implementation. Section~\ref{sec:accuracy62} establishes the investment primitives and the original-optimum benchmark. Sections~\ref{sec:query63} and \ref{sec:gated64} develop the query implementation; Section~\ref{sec:localization65} identifies its prediction and resource conditions. Section~\ref{sec:decisions49} states the economic inference contract, and Section~\ref{sec:study65} reports the complete service evidence. The technical supplement supplies independent algebraic checks, numerical contracts and complete comparison tables. Exact copies of the earlier article, supplement and broader development accompany this paper, with a label-level content map. Their diffusion, recursive-preference and game results retain their own hypotheses; the expected-cost query theorem is not silently extended to them.

\section{The economic operator and its policy account}\label{sec:core49}
Consider dates $t=0,\ldots,T$ and an invariant compact state domain $X\subset\R^d$. At date $t$, the action belongs to a nonempty compact set $A_t(x)$. The transition is $F_t(x,a,Z_{t+1})$, the current cost is $c_t(x,a)$, and the terminal cost is $g(x)$. Write $B_0=1$ and $B_t=\prod_{j<t}\beta_j$, with $\beta_j>0$. An implemented policy includes its observation rule, action selection and feasible numerical realization, not just its stored parameters.

For a continuation $v$, define
\begin{equation}\label{eq:operator49}
 Q_t(v;x,a)=c_t(x,a)+\beta_t\rho_t(v(F_t(x,a,Z))),\qquad
 \mathcal T_tv(x)=\inf_{a\in A_t(x)}Q_t(v;x,a).
\end{equation}
Throughout the deterministic comparison, the evaluations $\rho_t$ are monotone and cash invariant on a common domain containing all functions and constant translates used below. They are therefore nonexpansive in the supremum norm. Continuous objectives, effective compact feasible correspondences and measurable minimizers, or measurable arbitrarily near-minimizers with their error charged, are assumed. The stipulated conditional recursion gives Bellman comparison against all admissible adapted controls. In the experiments $\rho_t$ is conditional expectation, costs are bounded and innovations have their original continuous uniform law. The direct-cost inference in Section~\ref{sec:decisions49} is specific to expectation; it is not automatically an inference theorem for nonlinear recursive utility or games.

A Neural Bellman Operator is a construction and evaluation procedure for \eqref{eq:operator49}. It forms targets using its \emph{own} fitted $f_{t+1}$, returns a neural continuation and a feasible policy, and records an economic error allowance. It does not receive the optimal continuation as a training label. Arbitrary trained networks can be evaluated by the first result below. The subsequent constructive backend specifies how a network and its actor are obtained without assuming successful nonconvex optimization.

\subsection{Action-dependent errors and complete-policy loss}
Let $H(a)$ be a true current cost and $\widehat H(a)$ its continuation approximation. If $\widehat a$ is $\delta$-optimal for $\widehat H$ on the same feasible action set, then
\begin{equation}\label{eq:center49}
 H(\widehat a)-\inf_aH(a)\le\delta+
 \operatorname{osc}_{a}\{\widehat H(a)-H(a)\}.
\end{equation}
Indeed, compare $\widehat a$ with an arbitrarily near-optimal true action and cancel the approximate objectives. The infimum limit gives the result. Additive continuation levels cancel; action-dependent discrepancies do not. Under a change in future policies, the continuation must be transported to the new economic problem before this inequality licenses reuse. The full development edition retains the original quantitative transport, recursive-domain and economic-application results. No transport bound is presumed valid after an unaccounted change of constraints or preferences.

\begin{proposition}[One-sided policy sandwich]\label{prop:sandwich49}
Suppose, uniformly in state,
\[
 f_t-\mathcal T_tf_{t+1}\le u_t,\qquad
 \mathcal T_t^\pi f_{t+1}-f_t\le v_t,
 \qquad \ell_T\le f_T-g\le u_T.
\]
Then the actual feasible policy satisfies
\begin{equation}\label{eq:sandwich49}
 0\le J_0^\pi-V_0^*\le
 \sum_{t<T}B_t(u_t+v_t)+B_T(u_T-\ell_T).
\end{equation}
The bound is invariant to date-specific constant shifts of the fitted continuations, provided the associated residual accounts are shifted consistently.
\end{proposition}
\begin{proof}
Backward comparison and cash invariance give
\[
 V_t^*\ge f_t-\sum_{k=t}^{T-1}B_{t,k}u_k-B_{t,T}u_T,
 \quad
 J_t^\pi\le f_t+\sum_{k=t}^{T-1}B_{t,k}v_k-B_{t,T}\ell_T,
\]
where $B_{t,k}=\prod_{j=t}^{k-1}\beta_j$. The terminal claim is the stipulated enclosure; applying the Bellman or fixed-policy operator proves each inductive step. Subtraction gives \eqref{eq:sandwich49}. A shift $f_t\mapsto f_t+c_t^0$ adds $c_t^0-\beta_tc_{t+1}^0$ to $u_t$ and subtracts it from $v_t$. The terminal width is unchanged.
\end{proof}

A sharper valid residual account for the same continuation and the same deployed actor sharpens knowledge of the policy. It does not retrain that policy or change its realized cost. Conversely, subtracting two upper bounds in \eqref{eq:sandwich49} does not identify the difference in actual costs. These distinctions govern both the historical comparisons and the new experiments.

\section{Constructing and compiling a feasible neural policy}\label{sec:construction49}
Assume the primitive objective has, for feasible state--action pairs, state modulus $A_t^Q=C_{x,t}+\beta_tF_{x,t}L_{t+1}$ and action modulus $D_t=C_{a,t}+\beta_tF_{a,t}L_{t+1}$. Let
\[
 d_H(A_t(x),A_t(z))\le K_t^A\|x-z\|_1,
 \qquad L_t=A_t^Q+D_tK_t^A,\qquad L_T=L_g.
\]
The Hausdorff condition and the primitive moduli imply that $\mathcal T_tf_{t+1}$ is $L_t$-Lipschitz whenever $f_{t+1}$ is $L_{t+1}$-Lipschitz. All moduli and numerical integration procedures are supplied by the economic primitives, not fitted from successful outcomes.

At state nodes $x_{t,i}$ with covering radius $h_t$, form effective feasible $k_t$-nets. Evaluate their own-future objectives with certified error at most $e_t$. Minimize the computed labels and retain the corresponding \emph{original} action $a_{t,i}$. A Borel repair $\Gamma_{t,i}(x)$ is feasible at $x$ and satisfies
\begin{equation}\label{eq:repair49}
 \|\Gamma_{t,i}(x)-a_{t,i}\|_1\le K_t^A\|x-x_{t,i}\|_1+\zeta_t.
\end{equation}
The allowance $\zeta_t$ concerns displacement; feasibility itself may not fail. Define
\begin{equation}\label{eq:cone49}
 f_t(x)=\min_i\{y_{t,i}+L_t\|x-x_{t,i}\|_1\},\quad
 i_t(x)=\min\arg\min_i\{y_{t,i}+L_t\|x-x_{t,i}\|_1\}.
\end{equation}
Return $\pi_t(x)=\Gamma_{t,i_t(x)}(x)$. Terminal labels approximate $g$ with error $e_T$ and use the same cone construction with modulus $L_g$.

\begin{theorem}[Own-future construction with feasible witnesses]\label{thm:construct49}
The functions in \eqref{eq:cone49} have exact affine--ReLU realizations and Lipschitz moduli $L_t$ independent of label error. The returned policy is feasible and measurable. It satisfies
\begin{align}\label{eq:enclosures49}
 -e_t\le f_t-\mathcal T_tf_{t+1}&\le D_tk_t+e_t+2L_th_t,\\
 \mathcal T_t^\pi f_{t+1}-f_t&\le e_t+D_t\zeta_t,\notag
\end{align}
and hence
\begin{equation}\label{eq:constructionloss49}
 J_0^\pi-V_0^*\le
 \sum_{t<T}B_t(D_tk_t+2e_t+2L_th_t+D_t\zeta_t)
 +B_T(2e_T+2L_gh_T).
\end{equation}
With effective covers, repairs and certified queries at the prescribed accuracy, allocating each nonnegative term a share of a positive economic tolerance gives a terminating finite construction. A fixed numerical precision or a capped finite experimental ladder need not contain that allocation.
\end{theorem}
\begin{proof}
The selected node label lies between the node Bellman infimum minus $e_t$ and that infimum plus $D_tk_t+e_t$. The Lipschitz property of the infimum implies that every cone lies above it minus $e_t$. At a nearest node its cone exceeds the infimum by at most $D_tk_t+e_t+2L_th_t$. At the attaining original index, the two feasible pairs and \eqref{eq:repair49} give
\[
 Q_t(f_{t+1};x,\Gamma_{t,i}(x))
 \le y_{t,i}+e_t+L_t\|x-x_{t,i}\|_1+D_t\zeta_t.
\]
This proves \eqref{eq:enclosures49}. The terminal residual belongs to $[-e_T,e_T+2L_gh_T]$; Proposition~\ref{prop:sandwich49} gives the loss bound. A finite minimum of equally Lipschitz cones is equally Lipschitz. The identities $|z|=\operatorname{ReLU}(z)+\operatorname{ReLU}(-z)$ and $\min(u,v)=u-\operatorname{ReLU}(u-v)$ give the exact network. Finite lexicographic comparison and Borel repair give measurability. Effective refinement of the finitely many positively weighted terms yields the termination assertion.
\end{proof}

The Lipschitz-extension principle is classical \citep{mcshane1934}. A native min-plus envelope and its exact ReLU circuit, using the same labels, witnesses and ties, are \emph{one method with two encodings}. They return the same function and policy. Neither Theorem~\ref{thm:construct49} nor a faster representation proves a uniquely neural capability or a training-success probability for Adam or L-BFGS. The mathematical increment is the complete feasible, own-future construction and implementation account; the empirical question is its work and policy value relative to conventional generators.

\subsection{Original-owner compilation on a tensor cover}
On a tensor cover $\mathcal X=\prod_{j=1}^d\{x_{j,0},\ldots,x_{j,n_j}\}$, put $S=\prod_j(n_j+1)$. For arbitrary labels $y_i$, including inconsistent inexact labels, compute at each grid node $v$
\[
 (z_v,o_v)=\min_{\rm lex,i}(y_i+L\|v-x_i\|_1,i).
\]
The index $o_v$ is the original label owner, not the transformed node $v$.

\begin{theorem}[Continuous original-owner identity]\label{thm:compile49}
For a closed tensor cell $C$ with corners $\mathcal V(C)$, and $x\in C$,
\begin{align}\label{eq:compile49}
 f(x)&=\min_{v\in\mathcal V(C)}(z_v+L\|x-v\|_1),\\
 (f(x),i_*(x))&=\min_{\rm lex,v\in\mathcal V(C)}
 (y_{o_v}+L\|x-x_{o_v}\|_1,o_v).\notag
\end{align}
Coordinatewise forward and backward relaxations compute $(z,o)$ in $O(dS)$ additions and comparisons and $O(S)$ storage. After locating the cell, a query uses $O(d2^d)$ arithmetic operations. Exact rational comparisons preserve all ties, including $L=0$. Using the original owner's repaired action preserves the uncompiled policy, its cost and its certificate exactly.
\end{theorem}
\begin{proof}
The triangle inequality bounds $f(x)$ above by each corner expression. Choose the lexicographically first attaining site $i_*(x)$. In every coordinate choose the cell endpoint between $x_j$ and $x_{i_*,j}$. At this corner distances add exactly. An owner of strictly smaller transformed value would improve the objective at $x$; a tied smaller owner would contradict the original tie rule. Thus this corner has owner $i_*(x)$, proving both identities. Along a coordinate line, forward and backward relaxations respectively minimize over sites on either side, always comparing value--original-index pairs. Iterating over coordinates gives the separable distance transform and its stated work. Rational arithmetic and unchanged repair complete the claim.
\end{proof}

The nodal algorithm is the classical distance transform \citep{felzenszwalb2012}; preserving the original action through continuous off-grid queries is the policy obligation. If there are $J_{t,i}$ node actions and $M_t$ integration representatives, write $Q_t=M_t\sum_iJ_{t,i}$. Literal dense cone evaluation costs $O(\sum_t dS_{t+1}Q_t)$ in representation work. Compilation replaces this by $O(\sum_t[dS_{t+1}+Q_td2^d])$, plus cell location, primitive evaluation, integration, feasible repair, certificate and output. Rational bit lengths and state boxes crossing many cells incur additional costs. There is no dimension-free complexity assertion.

\subsection{Acquired states and conventional actors}
A $b$-bit coordinate sensor reports a dyadic cell and its midpoint $\widehat x$. On the unit cube, $\|x-\widehat x\|_1\le r(b)=d2^{-b-1}$. For the investment capacity $\kappa(x)$ below, clip the original action to the smallest capacity throughout that cell and round down to an action quantum $\Delta_t$. This is feasible at every state in the cell. Relative to an exact true-state repair, its displacement is at most $2K_t^Ar(b)+\Delta_t$. Transporting the witness cone from $\widehat x$ to $x$ adds at most $2L_tr(b)$, plus any certified selector objective error $\nu_t$. Thus the implemented-policy account adds
\begin{equation}\label{eq:acquire49}
 \sum_{t<T}B_t\{(2L_t+2D_tK_t^A)r(b)+D_t\Delta_t+\nu_t\}.
\end{equation}
No continuity of the discrete owner selector is assumed. At uncertain numerical comparisons the simulator retains all possible actors, rather than choosing a favorable one.

For conventional FVI the continuation is multilinear on its realized tensor grid. Its actual modulus is bounded by the largest coordinate edge slope. A nearest-node actor uses a feasible node minimizer and true-state repair. Let $h_t$ be the largest half-cell $\ell^1$ radius, $L_t^{\rm gr}$ the Bellman graph modulus, and $e_t$ the node query allowance. Its optimal residual upper bound is $D_tk_t+e_t+L_t^{\rm gr}h_t$. Its selected-policy upper bound is $e_t+E_t$, where
\[
 E_t=\sup_x\{y_{i(x)}+L_t^{\rm gr}\|x-x_{i(x)}\|_1-f_t(x)\}.
\]
Partitioning cells by nearest-node boundaries makes the expression separately affine in each coordinate. Its maximum is therefore at a partition vertex; the implementation encloses all such vertices. Both uniform and curvature-refined FVI receive this sharper one-sided account, not an artificially doubled baseline. Their acquired-state allowance is derived separately in \eqref{eq:fvisensor49}.

\section{Original-optimum accuracy in the investment economy}\label{sec:accuracy62}
The exact witness theorem answers an implementation question conditional on a continuation. An economist also needs to know how close the complete implemented policy is to the original optimum. We now establish that connection in the investment economy itself. The argument exploits a property not used by the earlier cell-range calculation: despite the bilinear transition, the Bellman objective remains jointly convex. This permits a second-order interpolation account and transfers verified nodal decisions to an implementable policy over the entire state domain. It does not assume that a fitted neural continuation is convex.

\paragraph{Relation to the preceding two-control query.}
The earlier query experiment used the fixed shared capacity $1/8$. The complete dynamic extension below uses the original state-dependent capacity $c(x)$, with the retained exposure columns, action-cost interaction and second shock. Its nesting claim concerns the original scalar economic model; it is not an assertion that a full dynamic policy was already present in the earlier isolated query.

\subsection{Primitives and Bellman regularity}
Let $d\ge2$ be even, let indices be cyclic, and set $\beta=15/16$. Write
\begin{align*}
 S(x)&=\frac2d\sum_j(x_j-5/8)^2+
       \frac1{4d}\sum_j(x_j-x_{j+1})^2+
       2\left(\frac12-\frac2d\sum_jx_j\right)_+^2,\\
 g(x)&=S(x)+\frac2d\sum_j(x_j-5/8)^2,\qquad
 c(x)=\frac18+\frac1{8d}\sum_jx_j.
\end{align*}
The original scalar control has $0\le a\le c(x)$, running cost $S(x)+a^2+4a^4$, and transition
\[
 F_j(x,a,z)=\frac1{16}+\frac{x_j}{2}+\frac{x_{j+1}}8+
       \frac{x_j(1-x_{j+1})}{16}+B_ja+s_jz,
\]
where $(B_j,s_j)=(1/2,1)$ at alternating coordinates and $(1/4,-1)$ at the others, and $z$ is uniform on $[-1/32,1/32]$.

For the retained two-control extension, $a_1,a_2\ge0$, $a_1+a_2\le c(x)$, the second exposure column interchanges $1/2$ and $1/4$, and the action cost is
\[
 R(a)=a_1^2+a_2^2+4(a_1^4+a_2^4)+a_1a_2/4.
\]
An independent common-direction uniform innovation of radius $1/64$ is added. This is the previously specified extension, nested at $a_2=0$ when that innovation is disabled. Both models preserve $[0,1]^d$. Let $m\in\{1,2\}$ denote the number of controls and define
\begin{align*}
 Q_t^*(x,a)&=S(x)+R(a)+\beta\mathbb E V_{t+1}^*(F(x,a,Z)),\\
 V_t^*(x)&=\min_{a\in A(x)}Q_t^*(x,a),\qquad V_T^*=g.
\end{align*}
All costs and laws in this display are the economic primitives, not a fitted surrogate.

\begin{theorem}[Regularity of the original Bellman problem]\label{thm:regularity62}
For either control specification and every finite horizon, $V_t^*$ is convex and continuously differentiable in the interior, with the corresponding one-sided extension on the cube. It is $G_t$-Lipschitz in the $\ell^1$ norm and $M_t$-semiconcave in the Euclidean norm, where valid constants are
\begin{align*}
 G_T&=10/d,&G_t&=\frac{1037}{128d}+\beta\frac{47}{64}G_{t+1},\\
 M_T&=26/d,&M_t&=\frac{22}{d}+\frac{21}{256d}
             +\beta\left(\frac9{16}M_{t+1}+\frac18G_{t+1}\right).
\end{align*}
In particular, $G_t\le27/d$ and $M_t\le54/d$. The function $Q_t^*$ is jointly strongly convex on its convex feasible graph, with state curvature at least $107/(128d)$ and action curvature at least $7/4$.

A valid semiconcavity constant along each coordinate line is
\begin{align*}
 D_T&=\frac9d+\frac{16}{d^2},\\
 D_t&=\frac5d+\frac{16}{d^2}+\frac{21}{256d^2}
 +\beta M_{t+1}\left(\frac{85}{256}+\frac{23}{256d}\right).
\end{align*}
Consequently, for the positive multilinear interpolant $I_n$ on a mesh of spacing $1/n$,
\begin{equation}\label{eq:interpolation62}
 0\le I_nV_t^*(x)-V_t^*(x)\le\kappa_t,
 \qquad \kappa_t=\frac{dD_t}{8n^2}.
\end{equation}
\end{theorem}

\begin{proof}
The state cost has gradient infinity norm at most $8/d$, Hessian norm at most $22/d$, and diagonal second derivatives at most $5/d+16/d^2$. For the terminal cost the corresponding bounds are $10/d$, $26/d$, and $9/d+16/d^2$. They follow by differentiating the quadratic terms and the continuously differentiable squared shortage. The bound on the cycle Laplacian is four. The shortage Hessian, where it exists, is $16\mathbf1\mathbf1^\top/d^2$; the same second-difference bound holds across its kink.

For feasible $(x,a)$ and $(y,b)$, put $w=x-y$. The difference between the transition at their convex combination and the corresponding combination of transitions is $\theta(1-\theta)w_jw_{j+1}/16$ in coordinate $j$. Convexity and the $G_{t+1}$ Lipschitz bound therefore give a convexity defect of at most
\[
 \frac{\beta G_{t+1}}{16}\theta(1-\theta)
                 \sum_j|w_jw_{j+1}|
 \le\frac{\beta G_{t+1}}{16}\theta(1-\theta)\|w\|_2^2.
\]
The state cost supplies curvature $4/d$, while the two-control action cost has Hessian eigenvalues at least $7/4$. Thus the state curvature of $Q_t^*$ is at least $4/d-\beta G_{t+1}/8$, which is at least $107/(128d)$ when $G_{t+1}\le27/d$. Partial minimization over the affine capacity graph preserves the state convexity inequality. The minimizing action is unique, including at a binding capacity constraint.

To propagate Lipschitz continuity, write $a=c(x)u$, where $u$ lies in the fixed unit simplex. Transfer the same $u$ from $x$ to $y$. Each action-cost partial derivative has absolute value at most $13/16$, and $\sum_j u_j\le1$. The action-cost transfer contributes at most $13\|x-y\|_1/(128d)$. The transition Jacobian at a fixed action has column sums at most $11/16$; the capacity transfer adds at most $3/64$. Comparing a minimizing normalized action in both directions gives the stated recursion for $G_t$. Its right side maps $27/d$ below itself.

For a fixed $u$, the transition as a function of $x$ has Jacobian norm at most $11/16+1/16=3/4$. Its second-order composition term is bounded by $G_{t+1}/8$. The action-cost Hessian has norm at most $21/4$, and the capacity gradient has squared norm $1/(64d)$. These facts give the recursion for $M_t$. The infimum over $u$ preserves this common semiconcavity bound: after subtracting the same quadratic, each fixed-$u$ function is concave, and an infimum of concave functions is concave. The right side maps $54/d$ below itself. Convexity together with semiconcavity yields the claimed differentiability and gradient regularity; this conclusion does not require a differentiable optimizing-action map.

Along one state coordinate, the fixed-$u$ transition has zero second derivative. The squared norm of its first derivative is at most $85/256+23/(256d)$: the two original nonzero entries are bounded by $9/16$ and $1/8$, and the capacity-induced entries are bounded by $1/(16d)$. Combining this with the diagonal stage-cost bound gives $D_t$. Finally, positive interpolation along one coordinate has error between zero and $D_t h^2/8$. Applying the coordinate interpolations successively and adding the bounds proves \eqref{eq:interpolation62}.
\end{proof}

The theorem is a structural statement about these economic primitives. It is not a claim that every controlled diffusion, recursive preference specification, or game in the complete development edition has a convex Bellman problem. Nor does it require a signed-weight neural fit to preserve convexity. Convex approximation of dynamic programs has an established literature; the additional work here is to verify the needed curvature for the original nonlinear capacity-constrained model and connect it to the actual NBO actor and complete error account.

\subsection{Certified nodal Bellman bounds}
Suppose nodal arrays satisfy $l_{t+1}\le V_{t+1}^*\le u_{t+1}$ at all grid vertices. Equation~\eqref{eq:interpolation62} gives the functions $I_nl_{t+1}-\kappa_{t+1}$ and $I_nu_{t+1}$ as valid lower and upper continuations. The lower function need not itself be convex; convexity is used for the unknown true value and for the comparison argument, not asserted of arbitrary endpoint arrays.

Divide each uniform innovation component into $q$ equal bins and use its conditional mean. Convexity gives a lower quadrature bound for the true continuation. Semiconcavity gives an upper error
\begin{equation}\label{eq:quadrature62}
 \zeta_{t+1}=\frac{M_{t+1}}{6q^2}
                      \sum_{r=1}^{s}\rho_r^2\|s_r\|_2^2,
\end{equation}
where $s=1$ or $2$, $\rho_1=1/32$, and $\rho_2=1/64$ in the extension. The expectation remains the original continuous law. At the terminal date $g$ is evaluated directly, so no terminal interpolation allowance is needed.

At a vertex $v$, let $\mathcal A_A(v)$ be the capacity simplex lattice with spacing $c(v)/A$, including all its faces. For $m=1,2$, a minimizer can be represented by neighboring lattice vertices whose coordinate variances sum to at most $m c(v)^2/(4A^2)$. A valid action smoothness constant is
\[
 H_t=\frac{21}{4}+\beta M_{t+1}\|B\|_2^2,
 \quad \|B\|_2^2=\frac{5d}{32}\ (m=1),\quad
 \|B\|_2^2=\frac{9d}{32}\ (m=2).
\]
Taylor's inequality and that barycentric representation imply
\begin{equation}\label{eq:action-cover62}
 0\le\min_{a\in\mathcal A_A(v)}Q_t^*(v,a)-V_t^*(v)
       \le\delta_t:=\frac{mH_t}{128A^2}.
\end{equation}
The linear Taylor terms cancel even when the optimum lies on a face. This is an error against the continuous feasible optimum, not only against the finite action menu.

Let $q_t^-(v,a)$ and $q_t^+(v,a)$ include all interpolation, quadrature and arithmetic errors and enclose $Q_t^*(v,a)$. Then
\begin{equation}\label{eq:nodal62}
 l_t(v)=\max\{0,\min_{a\in\mathcal A_A(v)}q_t^-(v,a)-\delta_t\},
 \qquad
 u_t(v)=\min_{a\in\mathcal A_A(v)}q_t^+(v,a)
\end{equation}
are valid nodal Bellman bounds. The zero truncation uses the nonnegative economic costs. Extra neural or conventional proposals may improve the upper side but do not remove the lower action-cover allowance. Positive interpolation, averaging and minimum maps are monotone and nonexpansive; the separate allowances in \eqref{eq:interpolation62}--\eqref{eq:action-cover62} are therefore propagated, rather than being discarded because an endpoint calculation contains the truth.

\subsection{From a fitted witness to the complete implemented policy}
Let $a_t(v)\in A(v)$ be the actual action returned at each vertex by any generator. It may be selected by the exact neural witness, a conventional fitted continuation, or their common verified fallback. Write $w_v(x)$ for the multilinear weights and set $\bar\pi_t(x)=\sum_vw_v(x)a_t(v)$. Affineness of capacity makes this action feasible. Let $q_t^{\mathrm{sel},+}(v)$ be a verified upper bound on its original-optimum continuation value $Q_t^*(v,a_t(v))$, using the same lower and upper Bellman reference for all generators.

\begin{theorem}[Actual-policy, whole-domain certificate]\label{thm:policy62}
Set
\[
 e_t=\max_v\{q_t^{\mathrm{sel},+}(v)-l_t(v)\}+\kappa_t.
\]
Suppose the implemented policy $\widehat\pi_t$ is feasible and differs from $\bar\pi_t$ by at most $b_t$ in action $\ell^1$ norm, uniformly over the state cube. Define
\[
 \Lambda_t=\frac{13}{16}+\beta\frac{3d}{8}G_{t+1},\qquad
 D_T^\pi=0,\quad D_t^\pi=e_t+\Lambda_tb_t+\beta D_{t+1}^\pi.
\]
Then, at every true state and date,
\[
 0\le J_t^{\widehat\pi}(x)-V_t^*(x)\le D_t^\pi.
\]
This statement is conditional on the actual recorded policy, not on success of its fitting algorithm. It applies equally to a pure fitted policy and to a policy augmented by verified common candidates.
\end{theorem}

\begin{proof}
Joint convexity, including the feasible graph, gives
\[
 Q_t^*(x,\bar\pi_t(x))\le\sum_vw_v(x)Q_t^*(v,a_t(v))
       \le I_nq_t^{\mathrm{sel},+}(x).
\]
The lower interpolation bound gives $V_t^*(x)\ge I_nl_t(x)-\kappa_t$. Subtraction proves the uniform optimal-advantage bound $e_t$. The partial action derivatives of $Q_t^*$ are bounded by $\Lambda_t$, so feasible implementation error adds at most $\Lambda_tb_t$. Adding the future policy loss and proceeding backwards proves the recursion. No convexity of $J^{\widehat\pi}$ is used.
\end{proof}

In the execution, observation is rounded down at forty bits. An acquired state $y\le x$ has no greater capacity than the true state. Each interpolated action is then rounded down at twenty bits using the lower endpoint of its arithmetic enclosure. If $L_{ij}$ bounds the derivative of action coordinate $i$ with respect to state coordinate $j$, and $\eta$ bounds native interpolation error, a valid implementation allowance is
\[
 b_t\le 2^{-40}\sum_{i,j}L_{ij}
           +m(2^{-20}+2\eta+\eta_{\mathrm{out}}).
\]
Here $L_{ij}$ is computed from the maximum adjacent nodal difference times $n$, and $\eta_{\mathrm{out}}$ covers endpoint subtraction. All these numbers are recorded outward. This new barycentric acquired actor is explicitly distinguished from the earlier discontinuous cell-constant actor; an old actor is not silently replaced during a purported fixed-policy revalidation.

\subsection{Arithmetic, work, and scope}
The native interpolation kernel forms nonnegative tensor weights and a weighted sum of at most $2^d$ stored binary64 endpoints. It is compiled without reassociation or floating-point contraction. For $K=(3d+4)2^d+32$, an explicit conservative absolute allowance is
\[
 \eta=\frac{K2^{-53}}{1-K2^{-53}}\max_v|z_v|+K2^{-1022}.
\]
The final term conservatively covers underflow. Input coordinates, affine indexing, finite endpoints and the binary64 format are checked. Roundoff expansion and the ordinary interval cost arithmetic remain separate from the analytic approximation bounds. Exact rational fixtures compare the native sum with its mathematical tensor interpolant.

With $N=(n+1)^d$ vertices, $K_A=A+1$ scalar actions or $K_A=(A+1)(A+2)/2$ two-control actions, the reference calculation uses $O(TNK_Aq^s2^d d)$ elementary work under the explicit tensor implementation. Stored reference values use $O(TN)$ words and a policy uses $O(TNm)$ words. Streaming batches avoid a state-pair table. Proposal fitting, exact rational witness work, actor transfer, serialization, independent cost inference and compilation are additional costs, not absorbed into this bound. The analytic errors are second order in state, action and innovation resolution; the state-count and tensor-evaluation factors still depend exponentially on dimension. The result supplies a matched original-optimum comparison without asserting that these costs have disappeared.

\section{Neural Bellman queries without a state lattice}\label{sec:query63}
The preceding construction makes numerical error accountable, but a complete tensor reference can absorb the benefit of a fitted continuation. We therefore change the computational organization, not the economic problem. A Neural Bellman selector proposes an action at the state actually acquired. A recursive interval query then bounds the original Bellman value at that state and checks the proposed action against the entire continuous feasible interval. No conventional policy is precomputed at all states. The controller is the resulting callable algorithm, not the unverified proposal alone.

The distinction matters for attribution. Learning may reduce the work needed to locate a useful action. It does not supply the inequalities used to certify that action. The adaptive conventional comparator receives the same inequalities and the same original-optimum target. The trained value-fit policies in the preceding study remain distinct from the action-query predictor introduced here.

\subsection{A curvature lower bound on an adaptive action partition}
Use the original scalar economy of Section~\ref{sec:accuracy62}, with $r$ dates remaining. Write $V_r$ for its optimal value and $Q_r(x,a)$ for the value of a first action followed by the optimal future. Let $G_r,M_r$ be the primitive constants in Theorem~\ref{thm:regularity62}, indexed by remaining horizon, and put
\[
 H_r=\frac{21}{4}+\beta M_{r-1}\frac{5d}{32},
 \qquad \Lambda_r=\frac{13}{16}+\beta G_{r-1}\frac{3d}{8}.
\]
Thus $H_r$ is an upper curvature bound in the scalar action, and $\Lambda_r$ bounds the absolute action derivative. Neither depends on a fitted model.

\begin{lemma}[Parabolic minorants]\label{lem:parabolic63}
Suppose $f$ is $H$-semiconcave on $[l,u]$, and $L_l\le f(l)$, $L_u\le f(u)$. For $s\in[0,1]$,
\[
 f(l+s(u-l))\ge (1-s)L_l+sL_u
                 -\frac H2(u-l)^2s(1-s)=:p(s).
\]
Consequently $\min_{[0,1]}p$ is a lower bound for the minimum of $f$ on the interval. With $K=H(u-l)^2/2$ and $D=L_u-L_l$, this lower bound is $L_l$ if $D\ge K$, $L_u$ if $D\le-K$, and otherwise
\[
 \frac{L_l+L_u}{2}-\frac K4-\frac{D^2}{4K}.
\]
The endpoints $L_l,L_u$ need not themselves belong to a convex interpolant. Increasing $K$ preserves the lower-bound property.
\end{lemma}
\begin{proof}
Apply concavity to $f(a)-Ha^2/2$ and substitute the two lower endpoints. The displayed quadratic has derivative $D-K+2Ks$. Minimizing it on the unit interval gives the three cases. Its coefficient of $K$ is $-s(1-s)$, which proves the last assertion.
\end{proof}

At an acquired state, retain an ordered partition of the feasible action interval and an outward value bracket at every partition point. The minimum of the parabolic minorants bounds the continuous-action infimum from below. The smallest upper endpoint supplies an actual feasible witness. A neural or polynomial proposal inserts another point into this partition; it cannot remove the endpoints or any part of the feasible interval. The running lower bound is the maximum of all previously valid global lower bounds, so new information cannot weaken the certificate. The upper bound is the minimum over every evaluated feasible witness.

\subsection{Recursive queries under the original continuous law}
For an endpoint action with $r>1$, divide the scalar innovation into $q_r$ equiprobable bins and query $V_{r-1}$ at each conditional-mean next state. If each child returns width at most $\eta_r$, the endpoint bracket has width at most
\begin{equation}\label{eq:querywidth63}
 \beta\eta_r+\frac{\beta M_{r-1}d}{6144q_r^2}+\omega_r.
\end{equation}
Here $\omega_r$ includes outward primitive arithmetic and the Lipschitz expansion from an interval-valued next state to its queried center. Convexity provides the lower conditional-mean inequality; semiconcavity and the exact conditional variance provide the upper allowance. The innovation remains continuous. There is no state-interpolation term because a child value is queried rather than read from a state lattice.

For a requested local gap $\delta$, the implementation sets $\eta_r\le\delta/(4\beta)$ and chooses a dyadic $q_r$ so that the second term in \eqref{eq:querywidth63} is at most $\delta/4$. A separate endpoint-width guard charges arithmetic and the small capacity-rounding remainder. Failure of that guard is a precision failure, not a successful return. The terminal query uses analytic integration instead of recursive quadrature.

\begin{theorem}[A state-grid-free NBO service]\label{thm:query63}
Consider the original scalar investment economy. At every recursive query, suppose the outward endpoint and capacity-transfer accounts give width at most $5\delta/8$. Let the lower bound be the running maximum described above. Refine an interval attaining the smallest current minorant whenever the global upper-minus-lower gap exceeds $\delta$. Each new point must lie between one fifth and four fifths of that interval. Arbitrary additional learned proposals may be inserted initially. Then:

(i) every returned interval $[L_r(x),U_r(x)]$ contains $V_r(x)$, and its returned feasible action $a_r(x)$ satisfies
\[
 Q_r(x,a_r(x))-V_r(x)\le U_r(x)-L_r(x)\le\delta;
\]
(ii) the real-arithmetic refinement terminates. With at most one interior learned proposal, a sufficient point budget is
\begin{equation}\label{eq:querycap63}
 6+\left\lceil5c(x)\sqrt{H_r/(3\delta)}\right\rceil;
\end{equation}
(iii) for the policy that calls this service at every subsequent date, the original-law loss satisfies
\begin{equation}\label{eq:querypolicy63}
 J_T^{\pi}(x)-V_T(x)
 \le\sum_{t=0}^{T-1}\beta^t
 \left[\delta_t+
 \left(8+\frac{11\beta dG_{T-t-1}}{16}
                  +dG_{T-t}\right)2^{-b}\right],
\end{equation}
when the state is acquired downward at $b$ bits per coordinate and every returned action is feasible at that acquired state. All action rounding must precede its endpoint evaluation. The inequalities hold uniformly in the true initial state, not only at the states used in an experiment.
\end{theorem}
\begin{proof}
At $r=1$, analytic endpoint enclosures and the minorant lemma give the assertion. Suppose it holds for $r-1$. Each child interval encloses the true optimal continuation. If the numerical next state is an interval, expanding the queried center value by $G_{r-1}$ times the maximum $\ell^1$ displacement encloses all its possible values. Conditional-mean lower bounds and the semiconcavity remainder then give valid endpoint brackets for the true $Q_r$. The minorants cover the entire action interval. A downward capacity endpoint can omit only an interval of length equal to the certified capacity-rounding difference; subtracting $\Lambda_r$ times that difference accounts for its possible improvement. The minimum upper endpoint is the cost of a feasible first action followed by the optimal future. This proves (i) by induction. Retaining earlier valid lower bounds preserves every conclusion.

For termination, an interval of length $h$ has minorant at least its smaller lower endpoint minus $H_rh^2/8$. Its smaller upper endpoint is an available global upper bound. With the endpoint-width guard, an interval of length at most
$h_\delta=\sqrt{3\delta/H_r}$ cannot be responsible for a gap exceeding $\delta$. Every refined parent is therefore longer than $h_\delta$. Its children are longer than $h_\delta/5$. Apart from the at most two initial fragments made by a prediction, disjoint terminal fragments consequently have total number at most $5c(x)/h_\delta$, up to endpoint counting. Equation~\eqref{eq:querycap63} is a conservative point budget. Bounded action-quantum rounding of a proposed quarter-to-three-quarter split preserves the one-fifth property whenever the quantum is smaller than $h_\delta/20$. Otherwise precision must be increased; a fixed-precision run is not allowed to ignore this requirement.

For (iii), let $\bar x$ be the acquired state. Its downward displacement is at most $d2^{-b}$ in $\ell^1$, and $c(\bar x)\le c(x)$ preserves feasibility. At a fixed action the stage cost is $8/d$-Lipschitz and the transition has column sums at most $11/16$. Thus replacing $\bar x$ by $x$ changes $Q_r$ by at most $(8+11\beta dG_{r-1}/16)2^{-b}$. The corresponding value change is at most $dG_r2^{-b}$. Add these two allowances to (i). Substituting the policy's own future yields the usual one-step loss recursion; backward induction gives \eqref{eq:querypolicy63}. No policy-value oracle and no population claim about fitting are used.
\end{proof}

The finite-precision contract is explicit. The executed service uses the width guard, a 512-point cap, and a 32-bit downward action lattice. A precision or cap failure withholds the floating certificate and invokes a separately counted exact-rational enumeration. This fallback constructs values only at the required descendants; it does not load a precomputed conventional policy. The rational fallback chooses dyadic action and quadrature counts until each cover or quadrature allowance is at most one quarter of the request, and requests child width at most one quarter divided by the discount. Provided $\delta>8\Lambda_r2^{-32}$ at every recursive level, its exact brackets have width at most three quarters of the request. This proves a finite return for every acquired state in the declared tasks, including states absent from the workload. Rational bit-operation cost is additional work and is not replaced by a unit-cost floating clock. The mathematical point budget is checked against the declared tolerance range; termination in an arbitrary floating implementation is not inferred merely from the real-arithmetic argument. Fitting convergence is unnecessary for validity: any finite prediction, and even its clipped replacement when numerically invalid, is only a proposed query.

\subsection{A certified analytic terminal kernel}
The last decision admits an especially strong common comparator. Let $y=F(x,0,0)$, $b_0=1/2-2\operatorname{mean}(y)$, and $\gamma_j=B_j$. Define
\[
 D(x)=\frac8d\sum_j\gamma_j(y_j-5/8)
   +\frac1{2d}\sum_j(\gamma_j-\gamma_{j+1})(y_j-y_{j+1}).
\]
For every even $d$, the exact terminal action derivative is
\begin{equation}\label{eq:cubic63}
 16a^3+\left(2+\frac{41\beta}{32}\right)a
 +\beta D(x)-3\beta(b_0-3a/4)_+.
\end{equation}
Each branch is a strictly increasing cubic. Its floating closed-form root supplies only a candidate. The independently evaluated outward derivative and exact-law objective certify that candidate; a safeguarded derivative bracket remains available if the candidate is inaccurate.

For completeness, if $f$ is $\mu$-strongly convex and an outward interval encloses $f'(a)$, the largest compatible projected derivative residual $R$ gives
\[
 f(a)-\min_{[0,c]}f\le R^2/(2\mu).
\]
At zero, a nonnegative derivative has residual zero; at capacity, a nonpositive derivative has residual zero. Otherwise the absolute derivative bound is sufficient. This follows by minimizing the strong-convexity supporting quadratic over the feasible interval and, conservatively, over the line. Here $\mu=2$. Equation~\eqref{eq:cubic63}, its $41/32$ and $9/4$ coefficients, analytic shock moments, and outward derivatives are independently checked against rational formulas. Every method uses this kernel; it is not credited to the neural predictor.

\subsection{Work, prediction, and the economic object}
Let $A_r(\delta)$ be a bound on action endpoints at a query. Excluding training, a scalar operation account obeys
\[
 W_1(\delta)=O(dA_1),\qquad
 W_r(\delta)\le C d A_r(\delta)+
 A_r(\delta)q_r W_{r-1}(\delta/(4\beta)).
\]
The constants $dG_r,dM_r$ are bounded independently of $d$, so the declared endpoint and quadrature budgets need not grow exponentially with the state dimension. Primitive evaluations still cost $O(d)$. This removes the state lattice, not the horizon or shock-tree cost. The executed vectorization batches a query's innovation descendants. Its largest live state batch is reported; it does not claim the smaller memory of a fully streamed implementation. Increasing the horizon or the shock dimension may still be expensive. The new execution is scalar-control; the earlier full two-control result is retained rather than relabeled as a query-service experiment.

A useful prediction can reduce the endpoint count or return its own action from the initial partition. An inaccurate prediction can instead add work. This is a prediction-dependent algorithm, not a theorem of neural dominance. Learned warm starts have their own literature, including \citet{sakaue2023warm}; their discrete-convex setting differs from the original-law Bellman intervals here. Convex control policies and learned optimization parameters are studied by \citet{agrawal2020control}. Our comparison identifies the extra proposal cost and the resulting query transcript rather than claiming priority for learning a warm start.

The purchase is now a complete callable controller with a declared error contract. Its training cost, every subsequent certified query, and its return-time cost belong to the same service. A reduction in queried state storage is economically relevant for repeated decisions only after these costs are counted. The original cost normalization is maintained. Workload seconds and normalized expected-cost loss are reported separately; a price per unit of computation can convert the former to the latter only as an explicitly specified decision-maker parameter.

\section{Certificate-gated neural queries}\label{sec:gated64}
A prediction of the optimal action and a prediction of a useful numerical query are different objects. Inserting a learned action before examining an already available certificate charges an extra continuation query even when the two boundary actions suffice. We instead make the demand for another query a mathematical decision and let a trained NBO selector determine its location only after that demand has been established. The underlying continuation recursion and the economic optimum are unchanged.

At a recursive state, start with the two feasible action endpoints. Compute the global lower bound from the parabolic minorants of Lemma~\ref{lem:parabolic63}, including the capacity remainder, and the smallest upper endpoint. Return immediately when their difference is at most the requested tolerance. Otherwise choose a partition interval with smallest current minorant. A learned action is used only if it lies in the central half of that interval; it then replaces the conventional proposed split. Outside that region the safeguarded parabolic split is used. The rounded split must pass the one-fifth interior check before evaluation. A prediction is evaluated at most once at each recursive state and is reused only as a proposed location. It supplies neither a value endpoint nor a lower bound.

\begin{theorem}[Gated prediction and complete-policy validity]\label{thm:gated64}
Under the assumptions of Theorem~\ref{thm:query63}, the gated service returns the same original-optimum and actual-policy guarantees for every supplied prediction, independently of fitting error. In real arithmetic, with endpoint-plus-capacity width at most $5\delta/8$, a sufficient number of endpoint evaluations at a state is
\[
 A_r(\delta,x)=4+\left\lceil5c(x)\sqrt{H_r/(3\delta)}\right\rceil.
\]
Each requested evaluation either reduces the unresolved action partition or terminates the query; no initial interior evaluation is mandatory. If the boundary certificate suffices, no predictor is evaluated at that state. If every prediction is declined at every recursive query, the value endpoints and implemented actions coincide with those of the conventional parabolic service using the same arithmetic and tie rules.
\end{theorem}
\begin{proof}
A routed action is merely another feasible partition point. The endpoint and minorant induction in Theorem~\ref{thm:query63} therefore applies without change. After each evaluation retain the smallest upper endpoint and the maximum of the old and newly valid global lower endpoints. Their difference cannot increase. The stopping test uses this difference, not a fitted value or a prediction error.

Put $h_\delta=\sqrt{3\delta/H_r}$. A partition interval of length at most $h_\delta$ has minimum minorant at least its smaller lower endpoint minus $3\delta/8$. The associated upper endpoint, with the stipulated width and capacity allowance, is within $\delta$ of that minorant. Such an interval cannot be responsible for an unresolved global gap. Every split parent therefore has length greater than $h_\delta$. The one-fifth condition makes each of its children longer than $h_\delta/5$. Starting with one interval, disjoint terminal fragments consequently number at most $1+5c(x)/h_\delta$. Counting endpoints and allowing for rounding of this numerical bound gives the displayed conservative budget. Floating failure is handled by the separately charged exact-rational procedure, not by this real-arithmetic count.

The first assertion about predictor evaluation follows from the order of the stopping test and the routing rule. For the last assertion, use backward induction on the remaining horizon, with the common analytic terminal kernel as base. At each nonterminal query, declined predictions leave the same conventional split, lower bound, witness and tie rule. Induction over partition refinements gives the same next query and returned endpoints; the induction over dates then gives the same implemented policy. Model-evaluation work can still differ, so transcript identity is not a zero-cost fitting claim.
\end{proof}

This theorem establishes robustness to the numerical content of a fitted prediction, not that learning is free or always useful. It also does not state that an arbitrary routed transcript has fewer queries than the adaptive conventional transcript. The executed comparison records both possibilities. Unlike a policy catalogue with a stored reference action at every state, neither service constructs a conventional state lattice before it can return a certified policy.

\subsection{State dimension, working memory, and reliability}
For the original scalar model, the primitive bounds $dG_r\le27$ and $dM_r\le54$ imply
\[
 H_r\le\frac{3369}{256},\qquad
 \Lambda_r\le\frac{1319}{128},\qquad
 \frac{\beta M_{r-1}d}{6144}\le\frac{135}{16384}.
\]
Thus the action and scalar-innovation budgets at fixed remaining horizon and tolerance are independent of the number of state coordinates. Every primitive evaluation still requires $O(d)$ arithmetic. A width-$w$ one-hidden-layer ReLU prediction on the recorded $d+5$ features requires $O(dw)$ work and stored parameters; a full quadratic prediction on these features uses $O(d^2)$ terms. Neither cost is included for free in the Bellman recursion.

Let $q_r(\delta)$ be the dyadic quadrature count and let $P_r(d)$ be the cost of one predictor evaluation. Apart from the separately identified exact-arithmetic recovery, a conservative work recursion is
\[
 W_r(d,\delta)\le C d A_r(\delta)+P_r(d)
   +A_r(\delta)q_r(\delta)
      W_{r-1}(d,\delta/(4\beta)).
\]
The predictor term is absent on a boundary-certified query. This expression removes the exponential state lattice, not the horizon or innovation-tree dependence. A depth-first scalar implementation can retain the current action partition and one innovation child at a time. The executed batched implementation instead reports its maximum state batch and process peak memory. A zero count of stored state-lattice nodes does not mean zero storage.

The deterministic certificate is valid for any supplied finite network weights. A failure to fit and return such weights is a separate service event. Our reliability estimand is the return count and worst measured complete work over the preregistered finite seed catalogue, with every repetition retained. Numerical workload paths are not optimizer replications. The real-arithmetic cap and the exact-rational precision condition establish algorithmic termination at admissible tolerances; recorded processor seconds are descriptive, not a hardware-uniform worst-case bound.

\subsection{An explicit economic purchase criterion}\label{sec:purchase64}
The decision maker buys a callable controller rather than a table of successful sample decisions. Fix a physical capital numeraire for the normalized state, a conversion $m>0$ dollars per unit of the model's loss index, a computation charge $p\ge0$ dollars per measured second, and a fixed deployment fee $f\ge0$ dollars. For $n$ uses under a specified initial-state law, the monetary objective is
\[
 \mathcal C(\pi)=nm\,\E[J_0^\pi(X_0)]+pW(\pi)+f.
\]
These are declared preference and service-price parameters, not estimated economic primitives. The policy theorem bounds the expected loss under the original continuous law for every such initial-state law. The finite dyadic workload below measures execution; its average cost is not substituted for that expectation.

If two policies have an independently valid expected-gain interval $[l,u]$ for $\E[J_0^{\pi_0}-J_0^{\pi_1}]$, and their incremental measured service charge is $\Delta=p(W_1-W_0)+f_1-f_0$, then the monetary net gain lies in
\[
 [nm l-\Delta,\;nm u-\Delta].
\]
This follows by multiplication by the positive numeraire conversion and subtraction of the same service charge. A positive lower endpoint supports purchase; a negative upper endpoint rules it out at the stated prices. Uniform regret certificates alone do not supply a directional expected-gain interval. They instead support a worst-case accuracy requirement imposed before selecting the least expensive eligible service.

Finally, a forecast of amortized query work is distinguishable from a measured complete service. With a fixed model and exactly repeated query workload, if training requires $B$ operations and deployment requires $D_L$ rather than $D_C$ operations per workload, then learning uses fewer total operations after $k$ repetitions precisely when $D_L<D_C$ and $k>B/(D_C-D_L)$. If $D_L\ge D_C$ and $B\ge0$, no repetition of that unchanged workload produces an operation saving. This arithmetic criterion includes training and cannot be applied to unrelated future states without a new assumption. The executed catalogue reports its own full process clocks, not a favorable hypothetical amortization.

\section{When a predicted query completes the certificate}\label{sec:localization65}
The safeguarded theorem is insensitive to the quality of a supplied prediction. A separate question is when that prediction is sufficient to finish verification. The answer must concern both the selected action and the unresolved lower bound. Proximity to a minimizer alone does not make a coarse lower certificate accurate.

\begin{proposition}[Witness accuracy and certificate localization]\label{prop:localization65}
Let $f$ be a differentiable, $H$-semiconcave objective with $H\ge0$ on $[0,c]$, let $a^*$ minimize $f$, and suppose $|f'|\le\Lambda$. An evaluated proposal $\widehat a$ lies in $[0,\underline c]$, where $0<\underline c\le c$ and $c-\underline c\le\Delta$. Each endpoint interval is a valid enclosure of $f$ and has width at most $w$. Suppose the partition containing $\widehat a$ covers $[0,\underline c]$ and has maximum interval length $h$. Let $L$ be its parabolic lower certificate with the capacity allowance $\Lambda\Delta$, and let $U$ be the smallest upper endpoint, retaining earlier valid bounds as in the query algorithm. With $e=|\widehat a-a^*|$ and $\lambda=|f'(a^*)|$,
\begin{equation}\label{eq:localization65}
 U-L\le 2w+\Lambda\Delta+\frac{Hh^2}{8}
                         +\lambda e+\frac{He^2}{2}.
\end{equation}
For an interior minimizer $\lambda=0$. At a boundary the linear term cannot generally be removed. Thus the right side at most $\delta$ is a sufficient certificate-return condition, not an assumption that a trained prediction meets it.

If old and new valid endpoint records give $(L,U)$ and $(L',U')$ after one additional query, with the running maximum lower bound and minimum upper bound, then
\begin{equation}\label{eq:certificate-progress65}
 (U-L)-(U'-L')=(U-U')+(L'-L)\ge0.
\end{equation}
An added prediction may have zero certified progress. The same identities apply to a conventional query.
\end{proposition}
\begin{proof}
Every retained lower endpoint is at least its true endpoint value minus $w$, and hence at least $f(a^*)-w$. On an interval of length at most $h$, the parabolic minorant is therefore at least $f(a^*)-w-Hh^2/8$. Subtracting the omitted-capacity allowance gives
$L\ge f(a^*)-w-Hh^2/8-\Lambda\Delta$; retaining an older, larger valid lower bound only improves this inequality. The upper endpoint at $\widehat a$ is at most $f(\widehat a)+w$. Semiconcavity and the one-sided first-order condition give
$f(\widehat a)-f(a^*)\le\lambda e+He^2/2$. Subtract the lower bound to prove \eqref{eq:localization65}. At an interior optimum the derivative is zero. The example $f(a)=a$ on $[0,c]$ shows that the boundary linear term is needed. Finally, retaining bounds makes $U'\le U$ and $L'\ge L$; expansion proves \eqref{eq:certificate-progress65}.
\end{proof}

This proposition isolates what a fitted representation can contribute without crediting it with a conventional solution. It can provide a small-error witness and change which unresolved interval is refined. The global partition, endpoint arithmetic and capacity account still have to support that witness. The maximum-width sufficient condition is deliberately conservative; the implemented minimum-minorant rule can stop earlier when only a subset of the partition remains competitive.

\subsection{A complete accuracy contract and its resource dependence}
For $T$ dates and target $\epsilon$, write $S_T=\sum_{t<T}\beta^t$. The recorded service chooses a common local tolerance by rounding $\epsilon/(2S_T)$ to binary64 and verifies the final outward allowance. Under $dG_r\le27$, a uniform per-date acquisition allowance at forty bits is
\[
 \alpha_b=\left(8+\frac{297\beta}{16}+27\right)2^{-40}.
\]
The service certifies $(\delta+\alpha_b)S_T\le\epsilon$. Each descendant request is $\delta/(4\beta)$, rounded as in the actual code and included in its outward endpoint guard. At the smallest request in the two- and three-date catalogues, the declared thirty-two-bit action quantum satisfies the exact-recovery condition $\delta_r>8\Lambda_r2^{-32}$. The supplement checks this inequality with rationals for every horizon and target, rather than inferring it from a successful workload.

For fixed remaining horizon and accuracy, the bounds on action endpoints and scalar-innovation bins are independent of state dimension in this normalized economy. The factors $dG_r$ and $dM_r$ are bounded, while every transition and cost still reads $d$ coordinates. A width-$w_\theta$ one-layer ReLU predictor has $O(dw_\theta)$ parameters and evaluation work; a dense quadratic on the same features has $O(d^2)$ terms. These statements identify the absence of a $(n+1)^d$ stored state table. They do not count an arbitrary-precision rational operation as a unit-cost floating operation and do not suppress the product of descendant query counts as horizon grows.

The two fitted implementations and their conventional controls therefore share one original-optimum contract, but need not share a transcript, a returned policy or a runtime. The completed experiments charge their distinct transcripts. A finite-catalogue reliability statement concerns all specified seeds and repetitions, including fitting warnings. It is not a population statement about optimizer success.

\section{From returned policies to economic resource decisions}\label{sec:decisions49}
An all-state upper bound and a comparison of actual policy costs answer different economic questions. We now connect the policies returned by a construction frontier to direct cost inference and then to a choice of implementation resources. The construction data select the policies; independently generated evaluation paths assess them. This separation permits a valid comparison of first-crossing policies with different resolutions.

\subsection{Simultaneous inference after construction}
Let $\mathcal H$ denote all construction information: primitive inputs, numerical labels, fitted continuations, returned policies and complete certificate ladders. Conditional on $\mathcal H$, fix a finite catalogue of implemented-policy contrasts. Each implementation includes its observation rule, original witness or nearest-node selector, true-state feasible repair and action quantization. For a contrast $i$, let $D_i$ be the difference in expected discounted costs under a specified initial law. A common-innovation path score evaluates the two policies on their own endogenous trajectories. It is not the difference between two regret bounds.

Suppose the initial and innovation uniforms are partitioned into equal bins. Independent uniform bin indices are drawn, and interval simulation encloses the score for every continuous realization inside the drawn bins. All unresolved state, selector and repair branches remain in the enclosure. If $[A_i,B_i]$ is a deterministic support interval, clip the path endpoints to that support. Write $\underline m_i,\overline m_i$ for outward sample-mean bounds on the lower and upper endpoint samples and $v_{i,L},v_{i,U}$ for upper bounds on their unbiased sample variances. For $n\ge2$ define
\begin{equation}\label{eq:radius49}
 R_i(v)=\sqrt{2v\ell/n}+\frac{7(B_i-A_i)\ell}{3(n-1)},
 \qquad \ell\ge\log(4K/\alpha),
\end{equation}
where $K$ bounds the number of sampled contrasts.

\begin{theorem}[Construction-conditional economic intervals]\label{thm:conditional49}
Under the stated independent-bin model, conditional on any valid construction history $\mathcal H$, with probability at least $1-\alpha$ all catalogue contrasts satisfy
\begin{equation}\label{eq:direct49}
 D_i\in[l_i,u_i]:=
 [\max\{A_i,\underline m_i-R_i(v_{i,L})\},
  \min\{B_i,\overline m_i+R_i(v_{i,U})\}].
\end{equation}
The same coverage holds unconditionally. Subsequent selection among these covered contrasts, or evaluation at any deterministic implementation-price vector, needs no additional union bound. This conclusion applies only to policies and initial laws included in the catalogue; it does not create direct-cost observations for every unstudied tolerance or initial state.
\end{theorem}
\begin{proof}
For each drawn bin sequence, conditional integration over its continuous realizations places the exact path expectation between the interval endpoints. Uniform mixing over the equal bins recovers the original continuous law. The lower endpoint expectation is thus no larger than $D_i$, and the upper endpoint expectation is no smaller. Apply the bounded empirical-Bernstein inequality of \citet{maurer2009} to the endpoint samples and take the union bound, with outward arithmetic included in the mean, variance and radius. Conditional validity for every $\mathcal H$ implies unconditional validity by integration. On this one simultaneous event, all ensuing deterministic comparisons of the covered quantities hold together.
\end{proof}

In particular, selecting the first policy to meet a specified certificate is permitted when that selection uses construction records rather than the subsequent evaluation paths. Repeated construction clocks are not independent draws from a neural-optimizer population. The statistical randomness here belongs to the evaluation design. A fixed pseudorandom stream implements a reproducible experiment; it is not itself a proof of the independent-bin assumption.

\subsection{Replacement-fee identification}
Write $D=J^W-J^F$. If witness is installed, replacing it by FVI changes total cost by $k-D$, where $k\ge0$ is the replacement fee. If FVI is installed, the change is $k+D$. Thus the two directional gains before fees are $D$ and $-D$.

\begin{corollary}[The break-even fee frontier]\label{cor:fee49}
On an event where $D\in[l,u]$, the two directional gains belong to $[l,u]$ and $[-u,-l]$. The unknown absolute cost difference satisfies
\begin{equation}\label{eq:absolute49}
 |D|\in[\operatorname{dist}(0,[l,u]),\max\{|l|,|u|\}].
\end{equation}
For every $k\ge\max\{0,u,-l\}$ neither replacement strictly lowers expected cost. By contrast, $k\ge\max\{0,l,-u\}$ only removes a \emph{certified} strictly profitable replacement. The two thresholds generally differ. Strict inequalities in the first criterion certify that both replacements strictly increase cost.
\end{corollary}
\begin{proof}
Subtract the gain intervals from $k$ and compare the resulting interval endpoints with zero. Taking the image of $[l,u]$ under the absolute-value map gives \eqref{eq:absolute49}. The assertions hold simultaneously for every nonnegative fee on the original coverage event.
\end{proof}

The fee frontier is more informative than a conclusion at one generous charge. The earlier $1/64$ replacement decision is retained with its original interpretation and confidence account. The new analysis reports the two directional gain intervals and both thresholds instead of selecting a favorable fee after seeing the data. There is no welfare calibration hidden in the units: a decision maker must supply the economically relevant charge, or use the full frontier to see what remains unidentified.

\subsection{Observation costs on the same bit grid}
Let $\mathcal B=\{4,\ldots,16\}$ be the available bit precisions. For one fixed returned controller, write $J_b$ for the expected cost of its actual $b$-bit, robustly repaired and quantized implementation. Use $b_0=16$ as a common reference and obtain simultaneous intervals for $D_b=J_b-J_{b_0}$. The identity interval at $b=b_0$ is exactly $[0,0]$. Let $s_b=d b\sum_{t<T}B_t$ be the discounted number of coordinate-bits purchased; the net cost is $J_b+\lambda s_b$.

\begin{theorem}[Certified choice of actual net implementation cost]\label{thm:net49}
Suppose all $D_b\in[l_b,u_b]$ simultaneously. For any $\lambda\ge0$ define
\begin{align}\label{eq:net49}
 L(\lambda)&=\min_b\{l_b+\lambda s_b\},&
 U(\lambda)&=\min_b\{u_b+\lambda s_b\},\\
 \widehat b(\lambda)&=\min\arg\min_b\{u_b+\lambda s_b\}.\notag
\end{align}
Then, simultaneously for all $\lambda\ge0$,
\begin{equation}\label{eq:netregret49}
 0\le J_{\widehat b}+\lambda s_{\widehat b}
       -\min_{b\in\mathcal B}(J_b+\lambda s_b)
 \le U(\lambda)-L(\lambda).
\end{equation}
The lower and upper functions are finite piecewise-affine envelopes. All changes in their active lines occur at nonnegative intersections of catalogue lines, with exact comparisons deciding ties. The smallest selected precision $\widehat b(\lambda)$ is weakly decreasing in $\lambda$.
\end{theorem}
\begin{proof}
Subtract the common $J_{b_0}$. The selected implementation has net difference at most $U(\lambda)$, and the best catalogue implementation has net difference at least $L(\lambda)$. This proves the regret bound pointwise on the one simultaneous event and hence for every $\lambda$. A minimum of finitely many affine functions changes its active affine expression only at an intersection. For $\lambda_2>\lambda_1$, sum the two minimizing inequalities to obtain $(\lambda_2-\lambda_1)(s_{\widehat b(\lambda_2)}-s_{\widehat b(\lambda_1)})\le0$. Since $s_b$ increases strictly with $b$, the precision comparative static follows.
\end{proof}

This result makes a statement about \emph{actual} net cost within the declared bit catalogue. It differs from minimizing a uniform all-state loss bound. Both criteria are useful: the latter gives distribution-free protection relative to the full-information optimum, while \eqref{eq:netregret49} uses a specified initial law to compare feasible implementations of one returned policy. Neither criterion gives an optimum over arbitrary sensor technologies or all policies.

For a fair comparison the conventional actor receives its own acquisition allowance. Suppose its nearest-node policy has a certified selected-policy account at a reported state, obtained from the graph modulus $L_t^{\mathrm{gr}}$, and its actual interpolated critic has modulus $L_t^f$. Transporting the same selected node action to the true state, then using cell-wide capacity repair, adds at most
\begin{equation}\label{eq:fvisensor49}
 (L_t^{\mathrm{gr}}+L_t^f+2D_tK_t^A)r(b)+D_t\Delta_t.
\end{equation}
Indeed, transport the feasible node pair directly to the true state, bound its node distance by the reported-state distance plus $r(b)$, and bound $f_t(\widehat x)-f_t(x)$ by $L_t^fr(b)$. Robust repair adds $D_t(2K_t^Ar(b)+\Delta_t)$. This proves \eqref{eq:fvisensor49} without assuming actor continuity. For a witness critic $L_t^f=L_t^{\mathrm{gr}}$, it reduces to the preceding witness allowance. The empirical comparison therefore evaluates the all-state-account minimizer, the interval upper-net-cost selector, and a descriptive midpoint selector on the \emph{same} integer bit and price grid. The midpoint selector is not called the true optimum.

\section{Complete services in the original investment economy}\label{sec:study65}
The numerical question is whether the organization of a learned query changes the work of delivering an accurate policy. We use the original scalar investment primitives, discount factor $15/16$, and continuous uniform innovations throughout. The returned object is the complete callable controller: at each subsequently acquired state it obtains a fresh certificate before implementing its action. A finite workload measures its execution. The uniform original-law theorem, rather than an extrapolation from that workload, supplies its accuracy guarantee.

\subsection{Design, controls and measured objects}
Two protocols were frozen before their respective executions. The first compared adaptive parabolic search, bisection, unconditional ReLU insertion and unconditional quadratic insertion. The second retained those controls and added certificate-gated routing for each fitted predictor. Its insertion and routing services independently reconstructed exactly the same fitted model for a matched seed, task and repetition. The initial and innovation workloads were also identical within each task. This design separates query organization from a change in model parameters or evaluated states.

The task catalogue is $(d,T,N)=(2,2,128),(8,2,128),(16,2,128),(8,3,32)$, with original-optimum targets $1/32$ and $1/128$. Here $N$ is the number of workload trajectories. Each learned method uses both specified fitting seeds, and every service has two isolated-process repetitions. The first protocol therefore contains 96 services and the second 160. The fitting seeds are 6301 and 6302 in the first catalogue and 6401 and 6402 in the second. No successful seed, task, target or repetition is selected for publication.

Each fitted service pays for its own 96 training states, nested Bellman-generated labels, feature construction and fit. The ReLU predictor has one hidden layer of width 24 and uses a capped L-BFGS fit. The quadratic predictor uses the same features, their second-order products and regularization. The permitted training iterations and function evaluations are fixed in the deposited protocols. At two dates the training account records 6,912 Bellman action queries; at three dates it records 155,904. These costs are not assigned to a previously solved conventional state lattice and are not omitted when a boundary certificate avoids a later prediction.

The adaptive comparator chooses an unresolved interval using its parabolic minorant and a safeguarded proposed split. Bisection supplies an independent refinement rule. Both choose their own action partitions and stopping stages; all methods use the same analytic terminal kernel, original-law child quadrature and precision guard. A fitted model cannot relax these requirements. These are structure-exploiting conventional controls for the scalar action problem, not a sparse-grid package relabeled as a neural method.

Every process is pinned to one CPU, with one numerical-library thread. The parent clock starts before launching the child and stops after process exit and join. It includes imports, training labels, fitting, all deployment queries, exact workload-cost reconstruction, serialized arrays, durable summaries and the final clock record. Processor frequency is uncontrolled. We report individual repetitions and machine-independent query counts alongside clocks; small clock differences are not treated as statistically identified runtime effects.

\subsection{Complete time and query attribution}
Table~\ref{tab:clocks65} reports the second catalogue's full cold-start process times. A conventional method has the smallest median in every task--target cell: bisection in three cells and adaptive parabolic search in five. None of the learned services has a smaller median complete time. The repeated reconstruction of a fitted model is part of this cold-start estimand. A different reuse contract would require a different work account, not deletion of the recorded fitting charge.

\begin{table}[htbp]
\centering\caption{Complete cold-start service times}\label{tab:clocks65}
{\small\setlength{\tabcolsep}{4pt}
\begin{tabular}{lrrrrrrrr}
\toprule
$d$ & $T$ & $\epsilon$ & A & B & RI & RG & QI & QG \\
\midrule
2 & 2 & 1/32 & 0.256 & 0.252 & 0.650 & 0.657 & 0.614 & 0.614 \\
2 & 2 & 1/128 & 0.270 & 0.282 & 0.655 & 0.656 & 0.624 & 0.628 \\
8 & 2 & 1/32 & 0.380 & 0.360 & 0.770 & 0.771 & 0.741 & 0.739 \\
8 & 2 & 1/128 & 0.416 & 0.441 & 0.806 & 0.806 & 0.754 & 0.773 \\
16 & 2 & 1/32 & 0.536 & 0.496 & 0.935 & 0.911 & 0.891 & 0.856 \\
16 & 2 & 1/128 & 0.552 & 0.622 & 0.978 & 0.987 & 0.913 & 0.935 \\
8 & 3 & 1/32 & 0.984 & 1.031 & 1.962 & 1.813 & 1.840 & 1.802 \\
8 & 3 & 1/128 & 1.450 & 1.952 & 2.435 & 2.481 & 2.323 & 2.285 \\
\bottomrule
\end{tabular}}
\par\smallskip{\footnotesize A: adaptive parabolic search; B: bisection; RI and RG: ReLU insertion and gating; QI and QG: quadratic insertion and gating. Seconds are parent-observed complete-process times, including imports, training, queries, exact workload costing, durable output and process join. Each conventional median uses two repetitions; each learned median uses two seeds and two repetitions. Single-host descriptive clocks are not significance tests. All individual observations are retained in the supplement.}
\end{table}

\begin{table}[htbp]
\centering\caption{Deployment Bellman query counts on identical workloads}\label{tab:queries65}
{\small\setlength{\tabcolsep}{4pt}
\begin{tabular}{rrrrrrrrr}
\toprule
$d$ & $T$ & $\epsilon$ & A & B & RI & RG & QI & QG \\
\midrule
2 & 2 & 1/32 & 1,298 & 1,310 & 1,427 & 1,298 & 1,421 & 1,295 \\
2 & 2 & 1/128 & 2,508 & 2,658 & 2,618 & 2,433 & 2,623 & 2,468 \\
8 & 2 & 1/32 & 1,280 & 1,325 & 1,472 & 1,278.5 & 1,461.5 & 1,278.5 \\
8 & 2 & 1/128 & 2,373 & 2,648 & 2,613 & 2,375.5 & 2,588 & 2,365.5 \\
16 & 2 & 1/32 & 1,286 & 1,328 & 1,455.5 & 1,281.5 & 1,457 & 1,283 \\
16 & 2 & 1/128 & 2,318 & 2,633 & 2,578 & 2,358 & 2,603 & 2,360.5 \\
8 & 3 & 1/32 & 7,042 & 8,112 & 9,014 & 6,947.5 & 8,604 & 7,038.5 \\
8 & 3 & 1/128 & 31,015 & 39,768 & 37,519 & 31,074 & 35,295.5 & 30,898.5 \\
\bottomrule
\end{tabular}}
\par\smallskip{\footnotesize Counts exclude training and are not a complete work measure. They include the recursive descendants and the common terminal kernel. Within a task and tolerance the workload is fixed across methods. Half-integer medians arise from two different seed transcripts. Insertion and gating share identical fitted parameters for each matched seed and repetition.}
\end{table}


The deployment counts in Table~\ref{tab:queries65} answer a narrower question. Holding each fitted model fixed, routing avoids the unconditional interior insertion and instead uses a prediction only when a certificate requires another query. There are sixteen distinct task--target--seed comparisons for each predictor. Routing reduces deployment Bellman queries in all sixteen ReLU comparisons and all sixteen quadratic comparisons. The repeated traces are identical, so repetitions do not turn these 32 matched comparisons into 64 independent numerical outcomes.

Against adaptive parabolic search the result is mixed. The ReLU route uses fewer queries in nine comparisons, the same number in three, and more in four. The quadratic route uses fewer in eleven, the same in two, and more in three. In the two-state, two-date, $1/128$ cell, the median ReLU deployment count falls from 2,508 for adaptive search to 2,433 for routing. In the sixteen-state counterpart it rises from 2,318 to 2,358. The three-date, tighter-target ReLU median also exceeds the adaptive count, despite improving substantially on unconditional insertion. The gate does not claim a pathwise query advantage over a strong conventional split rule.

\begin{table}[htbp]
\centering\caption{Matched query differences: gated minus comparator}\label{tab:matched65}
{\small\setlength{\tabcolsep}{4pt}
\begin{tabular}{llrrr}
\toprule
Predictor & Comparator & Lower & Equal & Higher \\
\midrule
Relu & Insertion & 16 & 0 & 0 \\
Relu & Adaptive & 9 & 3 & 4 \\
Quadratic & Insertion & 16 & 0 & 0 \\
Quadratic & Adaptive & 11 & 2 & 3 \\
\bottomrule
\end{tabular}}
\par\smallskip{\footnotesize One row comparison per task, target and fitting seed: sixteen for each predictor and comparator. Repetitions have identical models and numerical transcripts and do not double this count. These are query-count classifications, not signs of expected policy-cost differences.}
\end{table}


These observations identify an incremental contribution of query organization, not a full economic dominance result. A prediction may reduce the upper endpoint, improve the lower certificate through localization, or do neither. Proposition~\ref{prop:localization65} describes the sufficient mathematical conditions; the matched catalogue measures their algorithmic consequence without giving the predictor credit for the certificate's common components. Fitting and model evaluation can exceed the value of the saved queries. Here the complete clocks show that they do.

\subsection{Accuracy, dimension and finite-catalogue reliability}
All 256 declared services returned their required local certificates and uniform policy allowance. The second catalogue contains twenty services at each task--target cell. Its strongest state-dimensional execution is $d=16$, $T=2$, with target $1/128$ and no stored state-lattice nodes. This is different from the earlier eight-dimensional tensor construction, which attained only the coarse target $1/2$. The economic primitives and original optimum are unchanged; the computation now queries descendants instead of filling a state grid.

\begin{table}[htbp]
\centering\caption{Accuracy, memory and finite-catalogue return}\label{tab:reliability65}
{\small\setlength{\tabcolsep}{4pt}
\begin{tabular}{rrrrrrrr}
\toprule
$d$ & $T$ & Target & Returns & Gap bound & Batch & RSS (KiB) & Max (s) \\
\midrule
2 & 2 & 1/32 & 20 & 0.01562501 & 256 & 94520 & 0.945 \\
2 & 2 & 1/128 & 20 & 0.00390626 & 512 & 94496 & 0.670 \\
8 & 2 & 1/32 & 20 & 0.01562501 & 256 & 94748 & 0.792 \\
8 & 2 & 1/128 & 20 & 0.00390626 & 512 & 95160 & 0.824 \\
16 & 2 & 1/32 & 20 & 0.01562501 & 256 & 96040 & 0.953 \\
16 & 2 & 1/128 & 20 & 0.00390626 & 512 & 96212 & 1.006 \\
8 & 3 & 1/32 & 20 & 0.01562501 & 512 & 95704 & 1.964 \\
8 & 3 & 1/128 & 20 & 0.00390626 & 2048 & 98788 & 2.583 \\
\bottomrule
\end{tabular}}
\par\smallskip{\footnotesize Every cell has twenty isolated-process returns: two conventional methods with two repetitions and four learned methods with two seeds and two repetitions. Gap bounds are rounded upward for display; exact binary64 values are retained in the audit and checked against the rational target. RSS is the largest recorded process peak, not only model storage. No service built a stored state lattice or invoked exact-rational recovery.}
\end{table}


The uniform bound is approximately half the stated target plus the explicit acquisition allowance, because the fixed design reserves half of the target when selecting the per-date local request. Every exact recorded bound, not its printed decimal, is checked against its rational target. The proof covers true initial states outside the workload. The executed workload alone would not establish that assertion.

Removing a state lattice does not make the problem dimensionless. Each original cost and transition reads all coordinates, the dense quadratic model has a growing feature array, and the batched continuation recursion retains innovation descendants. Across both catalogues, the largest live batch is 2,048 states and the largest process peak is 98,860 KiB. The three-date cells reveal the continuation-tree cost; the tight adaptive deployment count rises to 31,015 queries for only 32 trajectories. No tensor-free two-control or high-dimensional shock execution is inferred from these scalar-control, scalar-innovation results. The earlier complete two-control policy study remains in the preserved article under its own implementation.

Reliability is defined for the frozen finite catalogue. All services returned; the worst recorded complete process time is 2.5831 seconds when rounded upward to four decimals. There were 104 fitting warnings across the two catalogues, all retained. A fitting warning does not forge a returned Bellman certificate, but it also cannot be discarded from a training-reliability description. No floating query invoked the separately counted exact-rational recovery on these workloads. That observation does not eliminate the recovery procedure from the uniform implementation contract or assign zero cost to its possible use elsewhere.

\subsection{Policy cost and the purchase decision}
The workload starts from fixed twenty-bit dyadic states and uses fixed dyadic shocks, with eight named boundary initial states. Every true state transition and path cost is reconstructed with rational arithmetic. There are 26,624 path rows and 55,296 implemented decisions across the two catalogues. These are finite diagnostic workloads, not independent draws furnishing new confidence intervals under the continuous law. The rational workload mean is retained for exact reproduction and is not substituted for the expectation in Section~\ref{sec:purchase64}.

The preceding continuous-law policy comparisons consequently remain important and separate. All twenty pure fitted ReLU policies in the tensor study have strictly higher expected cost than the common policy; the twenty common-augmented comparisons are unresolved. The common policy has the smallest complete release work in all twenty attainable task--target cells. The query study neither overwrites those observations nor asserts a new expected-cost reversal on the basis of finite dyadic averages. A callable query controller, a pure fitted actor, and a common-augmented nodal policy are distinct objects.

The purchase criterion has two uses. A buyer imposing only a uniform accuracy constraint can choose the least expensive eligible observed service for the declared workload and hardware. This favors a conventional method in the current cold-start catalogue. A buyer comparing monetary expected gains must additionally supply a valid expected-gain interval and the numeraire conversion, computation price and fixed fee in Section~\ref{sec:purchase64}. Neither the current clocks nor two regret upper bounds provide that directional gain interval. The criterion is a specified decision problem, not an empirical calibration of an installation charge.

\subsection{Reproduction and preserved evidence}
The revision checks every service request, source freeze, model identity, trace, exact cost and complete-process receipt. An independently written rational implementation reconstructs the economic trajectories. Re-querying each distinct recorded observation reproduces the original numerical endpoints, actions, probe counts and operation ledger; byte-identical repetitions are separately checked by identity rather than counted as fresh simulations. This is more than checking a table, but less than a formal proof-checker verification of arbitrary interval code.

The present integration creates no new training run, clock observation or independent continuous-law cost sample. The complete individual service tables, including fitting warnings and repeated timings, appear in the technical supplement. Every earlier manuscript source and its labeled results remain available in exact retained copies and four compiled companions. The active reading path brings the constructive NBO, primitive theorem, query implementation and economic comparison together without deleting an unfavorable experiment or extending a theorem beyond its hypotheses.


\section{Conclusion}\label{sec:conclusion65}
Neural Bellman Operators connect continuation construction to feasible economic decisions. The original witness construction supplies an explicit own-future loss account. In the nonlinear investment economy, primitive regularity also supports a recursive, state-lattice-free implementation. The trained component proposes where to obtain another piece of Bellman information; the certificate determines whether that information is needed and whether the returned action is adequate.

The resulting policy has a uniform original-optimum account even though only a finite workload is executed. That distinction relies on the full continuous-action cover, original-law quadrature, acquisition allowance and finite recovery contract, not on extrapolating from successful sample states. The cost still grows with the depth of the continuation tree. Removing a tensor state table is not removing horizon, innovation or arithmetic costs.

The matched evidence makes the role of learning assessable. Certificate gating improves the organization of learned queries, but its online savings do not pay for training on the observed cold-start workloads. This does not invalidate the construction or change its target; it determines which services the stated decision maker would purchase. Claims about repeated use require explicit saved work and an amortization condition, and claims about expected policy gain require independent economic inference. Keeping those obligations separate preserves the substantive NBO objective and makes the next comparison reproducible.
\bibliographystyle{ecta-fullname}
\bibliography{references}
\end{document}
